% D42
% Auto-generated from markdown source
% origin: ai

\documentclass[12pt, a4paper]{article}
\usepackage[utf8]{inputenc}
\usepackage[T1]{fontenc}
\usepackage{amsmath, amssymb, amsfonts, amsthm}
\usepackage{textcomp}
\usepackage{geometry}
\usepackage{hyperref}
\usepackage{booktabs}
\usepackage{graphicx}
\usepackage{fancyhdr}
\usepackage{tcolorbox}
\usepackage{enumitem}

\geometry{margin=1in}
\pagestyle{fancy}
\fancyhf{}
\rhead{D42}
\lhead{\thepage}

\newtheorem{definition}{Definition}[section]
\newtheorem{theorem}{Theorem}[section]
\newtheorem{corollary}{Corollary}[section]
\newtheorem{proposition}{Proposition}[section]
\newtheorem{lemma}{Lemma}[section]
\newtheorem{conjecture}{Conjecture}[section]

\newcommand{\calI}{\mathcal{I}}
\newcommand{\calP}{\mathcal{P}}
\newcommand{\calE}{\mathcal{E}}
\newcommand{\calS}{\mathcal{S}}
\newcommand{\calH}{\mathcal{H}}
\newcommand{\calO}{\mathcal{O}}
\newcommand{\calF}{\mathcal{F}}
\newcommand{\calU}{\mathcal{U}}
\newcommand{\calV}{\mathcal{V}}
\newcommand{\calA}{\mathcal{A}}
\newcommand{\calG}{\mathcal{G}}
\newcommand{\calB}{\mathcal{B}}
\newcommand{\calC}{\mathcal{C}}
\newcommand{\calL}{\mathcal{L}}
\newcommand{\calM}{\mathcal{M}}
\newcommand{\calX}{\mathcal{X}}
\newcommand{\calY}{\mathcal{Y}}
\newcommand{\calW}{\mathcal{W}}
\newcommand{\calT}{\mathcal{T}}
\newcommand{\calZ}{\mathcal{Z}}
\newcommand{\calR}{\mathcal{R}}
\newcommand{\calQ}{\mathcal{Q}}
\newcommand{\calD}{\mathcal{D}}
\newcommand{\calN}{\mathcal{N}}
\newcommand{\calK}{\mathcal{K}}

\newcommand{\Tr}{\text{Tr}}
\newcommand{\E}{\mathbb{E}}
\newcommand{\diff}{\mathrm{d}}
\newcommand{\KL}{D_{\mathrm{KL}}}
\newcommand{\argmax}{\arg\max}
\newcommand{\argmin}{\arg\min}
\newcommand{\hbarF}{\hbar}

\newtcolorbox{thmbox}[1]{
  colback=blue!5!white, colframe=blue!75!black, fonttitle=\bfseries,
  title={#1}, breakable
}
\newtcolorbox{defbox}[1]{
  colback=green!5!white, colframe=green!75!black, fonttitle=\bfseries,
  title={#1}, breakable
}

\begin{document}

\title{D42}
\author{RALPH Loop}
\date{}
\maketitle

\begin{abstract}
This paper presents D42, formalizing the underlying principles using information-theoretic methods.

\textbf{Keywords}: D42, information theory, formalization.
\end{abstract}

\tableofcontents
\section{D42}




\subsection{0.}


\textbf{D42 }

1. ****small-worldscale-free
2. ****KOL
3. \textbf{ = }betweennesscloseness
4. \textbf{ = } =  =



\subsection{1.}


\subsubsection{1.1}


\textbf{ D42-1 \(G_I\)}
\(V = \{v_1, v_2, \ldots, v_n\}\) \(E_I\)
\[P(v_i \to v_j) = \frac{w_{ij}}{\sum_k w_{ik}}\]

\(w_{ij}\)  \(v_i\)  \(v_j\)

\textbf{ D42-2 \(\Phi_u\)}
\[\Phi_u = \sum_{v \in \text{ancestors}(u)} \frac{S_v}{d(u, v)^\alpha}\]


\begin{itemize}
\item \(S_v\) \(v\)
\item \(d(u, v)\) \(u\)  \(v\)
\item \(\alpha\) \(\alpha \in [0.5, 2.0]\)
\end{itemize}


\textbf{ D42-3 \(H_\Phi\)}
\[H_\Phi = -\sum_u \frac{\Phi_u}{\sum_v \Phi_v} \log \frac{\Phi_u}{\sum_v \Phi_v}\]




\subsubsection{1.2}


\textbf{ D42-4 \(E(u)\)}
\[E(u) = \beta_0 + \beta_1 \Phi_u + \beta_2 \Phi_u \cdot A_u + \epsilon_u\]

\(A_u\)  \(u\) /\(\beta_2 > 0\) ****/

\textbf{ D42-5 \(T_u\)}
\[T_u = \frac{E(u)}{\Phi_u} = \frac{\beta_0}{\Phi_u} + \beta_1 + \beta_2 A_u\]



\textbf{ D42-6 \(\mathcal{M}\)}
\[\mathcal{M} = 1 - \frac{\text{Gini}(E)}{H_\Phi}\]

\(\text{Gini}(E)\) \textbf{\(\mathcal{M}\) }


\subsubsection{1.3}


\textbf{D42-1 }
\(\alpha_c\) \(\alpha < \alpha_c\)

****
\[\alpha_c = \frac{d \ln \langle \Phi \rangle}{d \ln n}\]

\(n\) \(\langle \Phi \rangle\)

****
\begin{itemize}
\item \(\alpha < \alpha_c\)scale-free
\item \(\alpha > \alpha_c\)/
\end{itemize}



\subsubsection{1.4}


\textbf{ D42-7 \(R_{between}\)}
\[R_{between}(u) = \gamma \cdot B(u) \cdot \max_v \Phi_v\]

\(B(u)\)  \(u\) betweenness centrality\(\gamma\)

\textbf{D42-2 }
****

****
1.  \(x\)  \(s\)  \(t\)
2. \(s \to t\)  \(x\)
3. \(x\)
4. \(R_{between}(x) = \text{flow}(s \to t) \cdot \text{tax}(x)\)
5.  ****


\subsection{2.}


\subsubsection{ D42-1 -}


****
\[\text{Gini}(E) \geq \text{Gini}(\Phi) \cdot (1 + \beta_2 \cdot \bar{A})\]

\(\bar{A}\)

****


\subsubsection{ D42-2}


****
\[\frac{\partial T_u}{\partial \Phi_u} = -\frac{\beta_0}{\Phi_u^2} + \frac{\partial \beta_2 A_u}{\partial \Phi_u}\]

****
-
-
\begin{itemize}
\item \textbf{ \(\Phi_u^* = \sqrt{\beta_0 / (\partial \beta_2 A_u / \partial \Phi_u)}\)}
\end{itemize}



\subsubsection{ D42-3}


****
\(m\)  \(\tau_m\)
\[\max_{\tau_m} \int_0^\infty e^{-rt} \cdot \text{flow}(t) \cdot \tau_m \, dt\]

\(\text{flow}(t) = F_0 \cdot e^{-\lambda \tau_m}\)

****
\[\tau_m^* = \frac{1}{\lambda}, \quad \text{flow}^* = F_0 \cdot e^{-1/\lambda}\]

****


\subsubsection{ D42-4}


****
\(\alpha \to \alpha_c^-\)
\[\lim_{\alpha \to \alpha_c^-} \mathcal{M}(\alpha) = 0\]

****

****


\subsection{3.}


\subsubsection{ D42-A}


\textbf{ D42-1}
1.  Newman-Watts small-world  \(p\)
2.  \(S_v = 1 \forall v\)
3.  \(p\)  \(\Phi_u\)  \(\alpha\)
4. ****\(p\)  scale-free


\subsubsection{ D42-B-}


\textbf{ D42-2-}
1.  \(\Phi_u \sim \text{Pareto}(\alpha_\Phi)\)
2.  \(A_u \sim \text{LogNormal}(\mu_A, \sigma_A)\)
3.  \(E(u) = \beta_1 \Phi_u + \beta_2 \Phi_u A_u\)
4.  \(\text{Gini}(E)\)  \(\text{Gini}(\Phi)\)
5. ****\(\text{Gini}(E) > \text{Gini}(\Phi)\) \(\beta_2\)


\subsubsection{ D42-C}


\textbf{ D42-3}
1.  \(m\)
2.  \(\text{flow}(\tau) = F_0 e^{-\lambda \tau}\)
3.  \(\max_\tau \text{flow}(\tau) \cdot \tau\)
4. ****\(\tau^* = 1/\lambda\)


\subsubsection{ D42-D}


\textbf{ D42-4}
1.  \(\alpha\)
2.  \(\mathcal{M}\)  \(\text{Gini}(\Phi)\)
3. ****\(\mathcal{M}\)  \(\alpha \approx \alpha_c\)


\subsection{4.}


\subsubsection{4.1 D8}


\textbf{D42-D8 }
\begin{itemize}
\item D8  \textbf{**  D42  }**
\item ****-
\end{itemize}



\subsubsection{4.2 D37}


\textbf{D42-D37 }
\begin{itemize}
\item D37  \textbf{ Nash }  D42  ****
\item ****
\end{itemize}



\subsubsection{4.3 D40}


\textbf{D42-D40 }
\begin{itemize}
\item D40  \textbf{-}  D42  \textbf{-}
\item ****
\end{itemize}



\subsection{5.}


\subsubsection{5.1}


1. \textbf{D42-1 }
2. \textbf{D42-2 }
3. \textbf{D42-3 }
4. \textbf{D42-4 }

\subsubsection{5.2}


1. ****
2. ****
3. ****
4. ****


\subsection{6.}


**** \(G = (V, E)\) \(S: V \to \mathbb{R}^+\) \(\Phi_u = \sum_{v} S_v / d(u,v)^\alpha\) \(E(u) = \beta_1 \Phi_u + \beta_2 \Phi_u A_u\)

****
1. ****\(\alpha < \alpha_c\)
2. \textbf{-}\(\text{Gini}(E) \geq \text{Gini}(\Phi) \cdot (1 + \beta_2 \bar{A})\)
3. **** \(\tau^* = 1/\lambda\)
4. ****\(\mathcal{M}(\alpha) \to 0\)  \(\alpha \to \alpha_c^-\)

\textbf{*}***


*D42 *

\begin{thebibliography}{99}
\bibitem{ref1} Shannon, C.E. (1948). A Mathematical Theory of Communication.
\bibitem{ref2} Cover, T.M. (2006). Elements of Information Theory.
\end{thebibliography}

\end{document}